\section{Related work}
\label{sec:related}

Appendix~\ref{app:related} adds two-phase sampling, verification bias, raters without a gold
standard, uncertainty in rankings, design sensitivity and stratified allocation.

\paragraph{Whether a benchmark measures what it claims.}
BetterBench grades benchmarks against criteria for design, implementation and documentation
\citep{reuel2024betterbench}, a complementary framework asks which evaluative claims a piece of
evidence supports \citep{salaudeen2025measurement}, and retrieval tools find the items in a
benchmark that bear on a particular use \citep{diddee2026benchbrowser}. That work sits upstream of
ours: we take the benchmark and its gold answers as given and ask whether the published ranking is
determined by the pipeline's own outputs, which a valid benchmark does not guarantee.

\paragraph{Configuration moves scores, and extraction is part of the configuration.}
Evaluation results move with prompt formatting \citep{sclar2023quantifying}, with small changes
that reorder leaderboards \citep{alzahrani2024benchmarks}, and with the many settings of a widely
used evaluation harness, answer extraction among them \citep{biderman2024lessons}. The sentence
that the 2024 verifier looks for descends from the convention introduced with
Minerva \citep{lewkowycz2022minerva}, and learned extractors have followed \citep{yu2024xfinder}.
The closest study on the extraction side concerns multiple-choice evaluation, where extraction
methods can disagree with human judgement and produce misleading model comparisons
\citep{molfese2025right}. We study free-form mathematical answers and ask which pairwise
comparisons can be established when no extraction method is taken as correct. Benchmarks of verifier accuracy \citep{li2025verifybench} and
verifiers trained to handle the formats that defeat regular expressions
\citep{liu2025compassverifier} go the other way: rank the verifiers and adopt the best. A better
verifier shrinks the discordant cells but does not settle the responses on which verifiers still
disagree.

\paragraph{Label-efficient evaluation with a proxy.}
Prediction-powered inference combines many machine-generated labels with a small human-labelled
sample to produce valid confidence intervals \citep{angelopoulos2023ppi}, with variants for hybrid
evaluation of language models \citep{fisch2024stratified}, active labelling
\citep{zrnic2024active} and synthetic labels \citep{boyeau2025autoeval}. Two papers target
rankings: \citet{chatzi2024ranking} return rank sets with a coverage guarantee from
pairwise comparisons by a strong model calibrated against human comparisons, and
\citet{lee2026judge} correct the bias that a judge's imperfect sensitivity and specificity
introduce into a score. Correcting a proxy's bias requires its error characteristics to carry over
to the models being scored; when they are common to all of them,
Proposition~\ref{prop:oneverifier} applies and no pair can be reversed relative to another. We
therefore bound instead of correcting, with homogeneity assumed only where verifiers with different
extraction rules agree and carried as a parameter that we vary. The cost is fewer conclusions,
each conditional on a tolerance a reader may dispute.

\paragraph{Partial identification and design sensitivity.}
The bounds of Section~\ref{sec:setup} are partial-identification statements
\citep{manski2003partial}, closest to classical results on contaminated data
\citep{horowitz1995identification} and on binary variables misclassified at unknown rates
\citep{molinari2008partial}, with the models as sub-populations. Concurrent
work applies partial identification to the prevalence of a latent attribute reported by a panel of
language models, calibrated against validation labels \citep{chen2026partial}; Appendix~\ref{app:related}
states how our design differs.
